% ID: SEL-D0007-BC-E06
% ESTADO: validada. S008, 2026-11-07; tres pasadas del mismo autor y revisión matemática, didáctica y visual.
\textbf{Datos.}\par\nopagebreak[4]
Sean $E$ ``producto electrónico'' e $I$ ``compra por Internet''. Tenemos $P(E)=0.55$, $P(I)=0.72$ y $P(E\mid I)=0.64$. Esta última es condicionada, no conjunta.

\textbf{a) Intersección.}\par\nopagebreak[4]
\[
P(E\cap I)=P(E\mid I)P(I)=0.64(0.72)=\boxed{0.4608}.
\]

\textbf{b) Unión.}\par\nopagebreak[4]
Restamos la intersección para no contarla dos veces:
\[
P(E\cup I)=P(E)+P(I)-P(E\cap I)
=0.55+0.72-0.4608=\boxed{0.8092}.
\]

\textbf{c) No Internet sabiendo no electrónico.}\par\nopagebreak[4]
Por De Morgan, $E^c\cap I^c=(E\cup I)^c$ y su probabilidad es $1-0.8092=0.1908$. La condición tiene probabilidad $P(E^c)=0.45$:
\[
P(I^c\mid E^c)=\frac{P(I^c\cap E^c)}{P(E^c)}
=\frac{0.1908}{0.45}=\boxed{0.424}.
\]
Comprobación mediante regiones disjuntas:
\[\begin{array}{c|cccc}\text{Región}&E\cap I&E\cap I^c&E^c\cap I&E^c\cap I^c\\\hline P&0.4608&0.0892&0.2592&0.1908\end{array}\]
 Suman 1; las dos últimas suman 0.45 y su última parte representa el 42.4\,\% de esa condición.
